\documentclass[11pt]{article}
\usepackage[margin=0.92in]{geometry}
\usepackage{amsmath,amssymb,amsthm,mathtools}
\usepackage{graphicx}
\usepackage{microtype}
\usepackage{enumitem}
\usepackage{xcolor}
\usepackage[colorlinks=true,linkcolor=blue!55!black,citecolor=blue!55!black,urlcolor=blue!55!black]{hyperref}

\newtheorem{theorem}{Theorem}
\newtheorem{proposition}{Proposition}
\newtheorem{lemma}{Lemma}
\newtheorem{remark}{Remark}
\newcommand{\R}{\mathbb R}
\newcommand{\Pone}{\mathcal P_1(\R)}
\newcommand{\Ptwo}{\mathcal P_2(\R)}
\newcommand{\W}{\mathcal W}
\newcommand{\E}{\mathcal E_{\log}}
\newcommand{\dd}{\,\mathrm d}
\DeclareMathOperator{\tr}{tr}

\title{The Centered Fixed-Trace GUE Large-Deviation Principle\\
\large A full proof of the conjecture in arXiv:2508.13698}
\author{Discovery run \texttt{fa\_banach\_001} -- \texttt{agent\_lane\_18}}
\date{August 11, 2026}

\begin{document}
\maketitle

\begin{center}
\fbox{\parbox{0.92\textwidth}{\textbf{Status: candidate new full solution.}
The conjectured large-deviation principle is proved at speed $n^2$ by a
direct microcanonical logarithmic-gas argument and contraction through the
centering map. Independent expert review is recommended, especially for the
standard quantile-patch lower-bound lemma recorded in full below.}}
\end{center}

\section{Source conjecture}

Chaintron and Lacker ask on PDF page 35 of \cite{CL} for the large deviations
of the centered empirical eigenvalue law of the fixed-trace Gaussian Unitary
Ensemble (GUE). Their exact setup and proposed rate are reproduced below.

\begin{figure}[ht]
  \centering
  \includegraphics[width=0.98\textwidth]{figures/open_problem_crop.png}
  \caption{The fixed-trace GUE conjecture in \cite[p.~35]{CL}.}
\end{figure}

Let $\mathcal H_n$ be the real Hilbert space of $n\times n$ Hermitian
matrices with Hilbert--Schmidt norm. Let $M_n$ be uniform on
\[
  \{H\in\mathcal H_n:\tr(H^2)=n^2\}.
\]
Write $x_1^{(n)},\ldots,x_n^{(n)}$ for the eigenvalues of $M_n/\sqrt n$ and
put
\[
 L_n=\frac1n\sum_{i=1}^n\delta_{x_i^{(n)}},\qquad
 \mathsf C\mu=(T_{-m_1(\mu)})_{\#}\mu,\qquad T_a(x)=x+a.
\]
Thus $\mathsf C L_n$ is precisely the centered empirical measure displayed
in the source.

For a probability measure $\mu$, define the signed mean and quadratic moment
\[
 m_1(\mu)=\int x\,\mu(\dd x),\qquad
 m_2(\mu)=\int x^2\,\mu(\dd x),
 \qquad
 \E(\mu)=-\iint\log|x-y|\,\mu(\dd x)\mu(\dd y),
\]
with $\E=+\infty$ when the integral diverges positively. The source uses
\[
 I_{\log}(\mu)=\E(\mu)+\frac12m_2(\mu)-\frac34.
\]

\section{Result}

We use the standard convention that an LDP with speed $n^2$ and rate $J$
means
\[
 -\inf_{G}J\le \liminf_{n\to\infty}\frac1{n^2}\log\mathbb P(X_n\in G)
\]
for every open $G$, and the reverse upper bound for every closed set.

\begin{theorem}[Fixed-trace GUE conjecture]\label{thm:main}
The sequence $(\mathsf C L_n)$ satisfies an LDP in $(\Pone,\W_1)$ with
speed $n^2$ and good rate function
\[
 J(\mu)=
 \begin{cases}
 I_{\log}(\mu)+\dfrac12(1-m_2(\mu)),
   &m_1(\mu)=0,\ m_2(\mu)\le1,\\[4pt]
 +\infty,&\text{otherwise}.
 \end{cases}
\]
Equivalently, on its effective domain,
\[
 J(\mu)=\E(\mu)-\frac14.
\]
This is exactly the rate conjectured in \cite[p.~35]{CL}.
\end{theorem}

The proof has two independent pieces. First, diagonalization turns the
fixed-trace matrix ensemble into a logarithmic gas on a Euclidean sphere.
Second, centering is a continuous map in the topology requested by the
source, and its contraction fiber can be computed explicitly.

\section{The spherical logarithmic gas}

Put
\[
 \Sigma_n=\Big\{x\in\R^n:\sum_{i=1}^n x_i^2=n\Big\}
\]
and let $S_n$ denote its surface measure. Weyl's integration formula,
followed by the coarea formula and the change from the eigenvalues of $M_n$
to those of $M_n/\sqrt n$, gives the unordered joint eigenvalue law
\begin{equation}\label{eq:spherical-density}
 Q_n(\dd x)=\frac1{Z_n}
 \prod_{1\le i<j\le n}|x_i-x_j|^2\,S_n(\dd x),
 \qquad x\in\Sigma_n.
\end{equation}
All dimension-dependent constants are absorbed into $Z_n$. Ordering the
eigenvalues only contributes a factor $n!$, whose logarithm is $o(n^2)$.

\begin{proposition}[Microcanonical log-gas LDP]\label{prop:micro}
Under \eqref{eq:spherical-density}, the empirical measure
$L_n=n^{-1}\sum_i\delta_{x_i}$ satisfies an LDP in $(\Pone,\W_1)$ with
speed $n^2$ and good rate
\begin{equation}\label{eq:K}
 K(\nu)=
 \begin{cases}
 \E(\nu)-\dfrac14,&m_2(\nu)\le1,\\[3pt]
 +\infty,&m_2(\nu)>1.
 \end{cases}
\end{equation}
Moreover,
\begin{equation}\label{eq:partition}
 \lim_{n\to\infty}\frac1{n^2}\log Z_n=-\frac14.
\end{equation}
\end{proposition}

The domain $m_2\le1$, rather than $m_2=1$, is essential. At every finite
$n$ one has $m_2(L_n)=1$, but $\W_1$ convergence can lose quadratic moment.
The lost moment is carried by one eigenvalue of order $\sqrt n$.

\subsection{Compactness and lower semicontinuity}

Let
\[
 \mathcal D=\{\nu\in\Pone:m_2(\nu)\le1\}.
\]
This set is $\W_1$-compact. Indeed, the quadratic bound gives tightness and
\[
 \sup_{\nu\in\mathcal D}\int_{|x|>R}|x|\,\nu(\dd x)\le R^{-1},
\]
which gives uniform integrability of first moments. It is also closed by
Fatou's lemma.

The logarithmic energy is lower semicontinuous on $\mathcal D$. A convenient
way to see the point at infinity is the identity
\begin{align}\label{eq:positive-kernel}
 h(x,y)&=-\log|x-y|+\frac12\log(1+x^2)
                    +\frac12\log(1+y^2)\ge0,\\
 \E(\nu)&=\iint h(x,y)\,\nu(\dd x)\nu(\dd y)
          -\int\log(1+x^2)\,\nu(\dd x).
\end{align}
The inequality follows from
$|x-y|^2\le(1+x^2)(1+y^2)$. The first term is lower semicontinuous by
Portmanteau, while the second is continuous under $\W_1$ convergence because
$\log(1+x^2)=o(|x|)$. Hence $K$ is lower semicontinuous on the compact set
$\mathcal D$, and is a good rate once the bounds are proved.

Finally, the classical semicircle variational principle says
\begin{equation}\label{eq:variational}
 \inf_{\nu\in\mathcal D}\E(\nu)=\frac14,
\end{equation}
with equality at the centered variance-one semicircle law $\sigma$. For
completeness, this follows immediately from the standard inequality
$I_{\log}(\nu)\ge0$: if $m_2(\nu)\le1$, then
\[
 \E(\nu)\ge\frac34-\frac12m_2(\nu)\ge\frac14,
\]
and the source records $m_2(\sigma)=1$ and $\E(\sigma)=1/4$.

\subsection{Upper bound}

For $\varepsilon>0$, set
\[
 \Phi_\varepsilon(\nu)=
 \iint\log(|x-y|\vee\varepsilon)\,\nu(\dd x)\nu(\dd y).
\]
On $\mathcal D$ this functional is continuous in $\W_1$; its logarithmic
growth at infinity is uniformly integrable under the quadratic bound. For
$x\in\Sigma_n$,
\begin{equation}\label{eq:truncate-upper}
 2\sum_{i<j}\log|x_i-x_j|
 \le n^2\Phi_\varepsilon(L_n)-n\log\varepsilon.
\end{equation}
Also $\log S_n(\Sigma_n)=O(n)$ by the exact sphere-area formula. Therefore,
for a closed $F\subset\Pone$,
\begin{equation}\label{eq:unnormalized-upper}
 \limsup_{n\to\infty}\frac1{n^2}log
 \int_{\{L_n\in F\}}\prod_{i<j}|x_i-x_j|^2S_n(\dd x)
 \le \sup_{\nu\in F\cap\mathcal D}\Phi_\varepsilon(\nu).
\end{equation}
The decreasing limit of $\Phi_\varepsilon$ as $\varepsilon\downarrow0$ is
$-\E$. Compactness of $F\cap\mathcal D$ and the usual decreasing-maxima
lemma for upper semicontinuous functions give
\begin{equation}\label{eq:unnormalized-upper2}
 \limsup_{n\to\infty}\frac1{n^2}\log
 \int_{\{L_n\in F\}}\prod_{i<j}|x_i-x_j|^2S_n(\dd x)
 \le \sup_{\nu\in F\cap\mathcal D}[-\E(\nu)].
\end{equation}
Taking $F=\Pone$ and using \eqref{eq:variational} yields
\[
 \limsup n^{-2}\log Z_n\le-\frac14.
\]

\subsection{Lower bound and the escaped eigenvalue}

We give the construction because it is the only noncanonical feature of the
argument. First suppose that $\nu$ has compact support, a bounded density,
finite logarithmic energy, and $m_2(\nu)<1$. Choose $n-1$ separated quantiles
$y_1<\cdots<y_{n-1}$ of $\nu$, and set
\begin{equation}\label{eq:outlier}
 b_n=\left(n-\sum_{i=1}^{n-1}y_i^2\right)^{1/2}.
\end{equation}
Then
\[
 b_n=\sqrt{n(1-m_2(\nu))}+o(\sqrt n).
\]
In particular, the empirical law of $(y_1,\ldots,y_{n-1},b_n)$ converges to
$\nu$ in $\W_1$: the last point has mass $1/n$ and contributes only
$O(n^{-1/2})$ to the first moment.

Take disjoint intervals of radius $c/n$ around the bulk quantiles, with $c$
fixed sufficiently small, and parametrize the positive cap of $\Sigma_n$ by
\begin{equation}\label{eq:cap}
 x_n=\left(n-\sum_{i=1}^{n-1}x_i^2\right)^{1/2}.
\end{equation}
The surface Jacobian is
\[
 \sqrt{1+|\nabla x_n|^2}=\frac{\sqrt n}{x_n},
\]
which stays bounded above and below on this patch. Its logarithmic surface
volume is $-O(n\log n)=o(n^2)$.

If the density of $\nu$ is bounded by $A$, then
$y_j-y_i\ge (j-i)/(An)$. Shrinking $c$ if necessary shows uniformly on the
patch that
\[
 \sum_{i<j<n}\log|x_i-x_j|
 \ge \sum_{i<j<n}\log|y_i-y_j|-O(n\log n).
\]
The quantile Riemann-sum lemma for the locally integrable logarithmic kernel
gives
\begin{equation}\label{eq:riemann}
 \lim_{n\to\infty}\frac2{n^2}
 \sum_{i<j<n}\log|y_i-y_j|
 =\iint\log|x-y|\,\nu(\dd x)\nu(\dd y)=-\E(\nu).
\end{equation}
All $n-1$ interactions with $x_n\asymp\sqrt n$ total only $O(n\log n)$.
Thus every $\W_1$-neighborhood $G$ of $\nu$ satisfies
\begin{equation}\label{eq:unnormalized-lower}
 \liminf_{n\to\infty}\frac1{n^2}\log
 \int_{\{L_n\in G\}}\prod_{i<j}|x_i-x_j|^2S_n(\dd x)
 \ge-\E(\nu).
\end{equation}

The regularity assumptions cost nothing. If $\E(\nu)<\infty$ and
$m_2(\nu)<1$, first truncate, then convolve with a smooth compactly supported
density, and if needed apply a dilation arbitrarily close to one. The
resulting bounded-density compactly supported measures converge in $\W_1$,
retain second moment strictly below one, and their logarithmic energies
converge to $\E(\nu)$. This is the standard approximation lemma for
logarithmic energy; it follows directly by first truncating the nonnegative
kernel $h$ in \eqref{eq:positive-kernel} and then using monotone convergence.
If $m_2(\nu)=1$, apply the preceding bound to
$(D_s)_{\#}\nu$, where $D_s(x)=sx$ and $s<1$, and let $s\uparrow1$. Indeed,
\begin{equation}\label{eq:dilation-energy}
 \E((D_s)_{\#}\nu)=\E(\nu)-\log s.
\end{equation}
The lower bound is vacuous when $\E(\nu)=+\infty$.

Applying \eqref{eq:unnormalized-lower} to $(D_s)_{\#}\sigma$ and letting
$s\uparrow1$ yields
\[
 \liminf n^{-2}\log Z_n\ge-\frac14.
\]
This proves \eqref{eq:partition}. Subtracting the normalizer from
\eqref{eq:unnormalized-upper2} and \eqref{eq:unnormalized-lower} proves the
LDP of Proposition \ref{prop:micro}. Since every $L_n$ lies in the fixed
compact set $\mathcal D$, exponential tightness is automatic.

\section{Centering and completion of the proof}

\begin{lemma}[Centering is continuous]\label{lem:center}
The map $\mathsf C:(\Pone,\W_1)\to(\Pone,\W_1)$ is $2$-Lipschitz.
\end{lemma}
\begin{proof}
For any coupling of $\mu$ and $\nu$,
$|m_1(\mu)-m_1(\nu)|\le\W_1(\mu,\nu)$. Translating the same coupling gives
\[
 \W_1(\mathsf C\mu,\mathsf C\nu)
 \le\W_1(\mu,\nu)+|m_1(\mu)-m_1(\nu)|
 \le2\W_1(\mu,\nu).
\]
\end{proof}

By Proposition \ref{prop:micro}, Lemma \ref{lem:center}, and the contraction
principle, $\mathsf C L_n$ has good rate
\begin{equation}\label{eq:contraction}
 J(\mu)=\inf\{K(\nu):\mathsf C\nu=\mu\}.
\end{equation}
If $m_1(\mu)\ne0$, the fiber is empty. If $m_1(\mu)=0$, its members are
exactly the translates $(T_a)_{\#}\mu$, $a\in\R$. Translation invariance of
the logarithmic energy and centering of $\mu$ give
\begin{equation}\label{eq:fiber}
 \E((T_a)_{\#}\mu)=\E(\mu),\qquad
 m_2((T_a)_{\#}\mu)=m_2(\mu)+a^2.
\end{equation}
The fiber contains an admissible member of $\mathcal D$ exactly when
$m_2(\mu)\le1$. In this case every admissible member has the same $K$-value,
so \eqref{eq:contraction} becomes
\[
 J(\mu)=\E(\mu)-\frac14
 =I_{\log}(\mu)+\frac12(1-m_2(\mu)).
\]
This proves Theorem \ref{thm:main}.

\section{Checks, consequences, and scope}

\begin{enumerate}[leftmargin=*]
\item \textbf{Speed.} The Vandermonde has $\binom n2$ interactions, so the
  speed is $n^2$. The final display on source page 35 has
  $\limsup_{m\to\infty}n^{-1}\log$; this appears to be an $m/n$ and speed
  typographical error. The LDP proved here gives $n^{-2}\log$.

\item \textbf{Why no quadratic penalty remains.} On the fixed sphere, the
  quadratic external-field term is constant. Algebraically,
  $I_{\log}+(1-m_2)/2=\E-1/4$. Geometrically, missing quadratic moment can
  escape through one coordinate without an $n^2$-scale interaction cost.

\item \textbf{Centering does not hide a discontinuity.} The first moment is
  continuous in $\W_1$, and Lemma \ref{lem:center} gives a quantitative
  continuity estimate. No $\W_2$ contraction is invoked.

\item \textbf{Concentration consequence.} Combining this LDP with the
  lower-semicontinuous extension of the strengthened Biane--Voiculescu
  inequality in \cite{CL} gives the tail bound proposed there, with corrected
  speed $n^2$.

\item \textbf{Scope.} The proof is for the complex Hermitian ($\beta=2$)
  ensemble stated in the source. The same argument formally replaces $\E$
  by $\frac\beta2\E$ for fixed-trace $\beta$-Hermite laws, but no such
  extension is claimed here.
\end{enumerate}

\section{Novelty search and review recommendation}

A bounded search on August 11, 2026 checked the exact conjecture, the arXiv
identifier, ``fixed trace GUE empirical measure large deviations'',
``fixed trace beta-Hermite LDP'', and the lost-second-moment rate. It found
the classical canonical Wigner empirical-measure LDP \cite{BG}, work on
fixed-trace beta-Hermite limiting density and edge behavior
(arXiv:0905.4255), and large deviations for traces/moments
(arXiv:1605.03894), but no source stating the $\W_1$ microcanonical LDP
\eqref{eq:K} or its centered contraction. The result is therefore labeled a
candidate new full proof rather than a literature answer.

The recommended expert audit should focus on the quantile-patch estimate
\eqref{eq:riemann}, the regularization step following
\eqref{eq:unnormalized-lower}, and whether a later fixed-trace ensemble paper
contains the same theorem under different terminology. The contraction
calculation itself is exact and elementary.

\begin{thebibliography}{9}
\bibitem{CL}
L.-P. Chaintron and D. Lacker,
\emph{Geodesic convexity and strengthened functional inequalities in
submanifolds of Wasserstein space}, arXiv:2508.13698 (2025).

\bibitem{BG}
G. Ben Arous and A. Guionnet,
\emph{Large deviations for Wigner's law and Voiculescu's non-commutative
entropy}, Probab. Theory Related Fields \textbf{108} (1997), 517--542.

\bibitem{BV}
P. Biane and D. Voiculescu,
\emph{A free probability analogue of the Wasserstein metric on the
trace-state space}, Geom. Funct. Anal. \textbf{11} (2001), 1125--1138.
\end{thebibliography}

\end{document}
